\documentclass{jsarticle}
\usepackage{amsmath}
\usepackage{amssymb}
\begin{document}
\title{放射性物質吸収体とRNA干渉としてのオイラーの公式}
\author{Masaaki Yamaguchi}
\date{}
\maketitle

  　オイラーの公式から、XY平面での曲平面をZ軸での虚軸として、
  ゼータ関数をZ軸にAssembile-D-braneのシリンダーを束ねると、
  DNAのオミクロン曲線へと２重螺旋構造を再現できて、
  素数の分布が双対としてのDNAの父親と母親から受け継いでいる、
  染色体として、この説明でのアミノ基の構造が、ヒッグス場とオイラーの定数
  との加群で作られるガンマ関数における大域的微分多様体として、
  表現できることを、
  $$ \int \Gamma(\gamma)^{'}dx_m = {d \over df}F + \int C dx_m $$
  $$ = e^f + e^{-f} \ge e^{f} - e^{-f} $$
  $$ e^{-\theta} = \cos(i x \log x) - i \sin(i x \log x) $$
  $$ e^{i \theta} = \cos \theta + i \sin \theta $$
  これらの式たちから、DNAの生体地図としての、染色体の構造式は、
  オイラーの公式と同型としての、
  $$ \int \Gamma(\gamma)^{'}dx_m  = 2(\cos (i x \log x) - i \sin (i x \log x))$$
  と表される。

  この式は、反重力が放射性物質から励起される熱エネルギーを
  、カーボンナノチューブとコバルト６０との
  複合体でもある、二量体としての結晶石が、吸収することができる、
  生体エネルギーの人体の保温と同じ原理で、反重力の放射性物質を吸収
  することがこの結晶石が、オイラーの定数の大域的微分多様体から
  導く金属の結晶構造の配置と同じ原理で表される式として、
  ヒッグス場とオイラーの定数の大域的微分多様体の式が、
  この二量体の結晶石の構造式として、求まる。

  これらの式たちは、次の理論で証明できる。
  数値は結果であり、アイデアから式、そして値をもとめることが、
  エレガントな解き方と言える。
  

    ベータ関数をガンマ関数へと渡すと、
   $$ \beta(p,q) = {{\Gamma(p)\Gamma(q)} \over {\Gamma(p + q)}} $$
   ここで、単体的微分と単体的積分を定義すると
   $$ t = \Gamma(x) $$
   $$ T = \int \Gamma(x) dx $$
   これを使うのに、ベータ関数を単体的微分へと渡すために、
   $$ \beta(p,q) = - \int {1 \over {t^2}}dt $$
   大域的微分変数へとするが、
   $$ T_m = \int \Gamma(x) dx_m $$
   この単体的積分を常微分へとやり直すと、
   $$ T^{'} = {t^{'} \over {{\log t}}}dt + C(\text{Cは積分定数}) $$
   これが、大域的微分多様体の微分変数と証明するためには、
   $$ dx_m = {1 \over {\log x}}dx, dx_m = {(\log x^{-1})}^{'}  $$
   これより、単体的微分が大域的部分積分へと置換できて、
   $$ T^{'} = \int \Gamma(\gamma)^{'} dx_m $$
   微分幾何へと大域的積分と大域的微分多様体が、加群分解の共変微分へと
   買い換えられる。
   $$ \bigoplus T^{\nabla}  = \int T dx_m, \delta(t) = t dx_m  $$


   これらをまとめると、ベータ関数を単体的関数として、
   これを常微分へと解を導くと、単体的微分と単体的積分
   へと定義できて、これより、大域的微分多様体と大域的積分多様体
   へとこの関数が存在していることを証明できるのを上の式たちから
   わかる。この単体的微分と単体的積分から微分幾何が書けることがわかる。

   この説明から単体的微分と単体的積分が結晶石の二量体としてと、
   RNA干渉としての、オイラーの公式をも、構造式として表せることを
   理論建てしている。

   アメリカの国防総省のUFOの解体から、リバースエンジニアリングした
   UFOもどきが、この航空機が飛行するのに、機体の中心部に
   放射性物質による結晶石をつかっているのが、宇宙人の発想の逆を言っている
   とわかり、早坂先生の反重力発生装置が解決しているのを、
   この反重力発生装置の熱エネルギーを発する放射性物質自体の
   航空機への放射性物質の被爆を防ぐのに、UFOは中心部へこの放射性物質
   を吸収する結晶石を置いているのと、その放射性物質の
   副次的エネルギーを、UFOの副次的動力源として、
   エネルギーのリサイクルをも、人間が核エネルギーの
   原子力発電所のリサイクルエネルギーと同じ発想で、
   UFOの動力源の補助をもこの放射性物質の結晶石がしているというのが、
   わかった次第でもあります。
   なぜ、このUFOもどきをつくっても少ししか上空しないかは、副次的エネルギー
   としての、主エネルギーに比べて、少ないという、この機知
   の仕組みで歴然としています。

   UFOが地球に頻繁に現れ始めたのが、原爆のあとと、ATR研究所が
   ブレインインターフェースを作り始めたら、宇宙人による人へのインプラント
   が頻繁に始まり、本当に理論と工学のからくりがわかりやすく、
   一人さんの持論を宇宙人までしている始末でもあります。

\end{document}
